\subsection*{الف}
فرض کنید $x(t)$ یک پالس مستطیلیِ واحد با دوره‌ی $T$ و پهنای $\tau$ (با $0<\tau<T$) باشد:
\begin{flalign*}
x(t)=\sum_{k=-\infty}^{\infty}\Big(u(t-kT)-u(t-kT-\tau)\Big), \qquad \omega_0=\frac{2\pi}{T}. &&
\end{flalign*}
ضرایب سری فوریه‌ی نمایی:
\begin{flalign*}
c_n=\frac{1}{T}\int_{0}^{T} x(t)e^{-jn\omega_0 t}\,dt
      =\frac{1}{T}\int_{0}^{\tau} e^{-jn\omega_0 t}\,dt. &&
\end{flalign*}
پس
\begin{flalign*}
c_0&=\frac{\tau}{T},\\
c_n&=\frac{1}{T}\,\frac{1-e^{-jn\omega_0\tau}}{jn\omega_0}\qquad (n\neq 0). &&
\end{flalign*}
مشتقِ $x(t)$ (به معنای تعمیم‌یافته) یک جمع از ضربه‌هاست:
\begin{flalign*}
x'(t)=\sum_{k=-\infty}^{\infty}\Big(\delta(t-kT)-\delta(t-kT-\tau)\Big). &&
\end{flalign*}
اگر $d_n$ ضرایب سری فوریه‌ی $x'(t)$ باشد، با استفاده از این که ضرایب سری فوریه‌ی
$\sum_k \delta(t-t_0-kT)$
برابر $\frac{1}{T}e^{-jn\omega_0 t_0}$ است، داریم:
\begin{flalign*}
d_n=\frac{1}{T}\Big(1-e^{-jn\omega_0\tau}\Big). &&
\end{flalign*}
اکنون بررسی رابطه:
\begin{flalign*}
jn\omega_0\,c_n
=jn\omega_0\left(\frac{1}{T}\frac{1-e^{-jn\omega_0\tau}}{jn\omega_0}\right)
=\frac{1}{T}\Big(1-e^{-jn\omega_0\tau}\Big)
=d_n \qquad (\forall n). &&
\end{flalign*}
پس در این مثال (با درنظرگرفتن مشتقِ ضربه‌ای)، رابطه‌ی
$d_n = j\frac{2\pi n}{T}c_n$
صحیح است.

\subsection*{ب}
فرض کنید $f(t)$ با دوره‌ی $T$ در نقاط $t_i$ (در یک تناوب) ناپیوستگی پرشی داشته باشد و
\begin{flalign*}
\Delta_i=f(t_i^+)-f(t_i^-). &&
\end{flalign*}
ضرایب سری فوریه‌ی $f(t)$ را $c_n$ بگیرید. اگر $g(t)=f'(t)$ مشتقِ معمولیِ $f$ در نواحیِ هموار باشد
(یعنی بدونِ ضربه‌ها)، آنگاه با جزءبه‌جزء کردن انتگرال و انتگرال‌گیری جزءبه‌جزء:
\begin{flalign*}
g_n&=\frac{1}{T}\int_{0}^{T} f'(t)e^{-jn\omega_0 t}\,dt \\
&=\frac{jn\omega_0}{T}\int_{0}^{T} f(t)e^{-jn\omega_0 t}\,dt
-\frac{1}{T}\sum_i \Big(f(t_i^+)-f(t_i^-)\Big)e^{-jn\omega_0 t_i} \\
&= jn\omega_0\,c_n-\frac{1}{T}\sum_i \Delta_i\,e^{-jn\omega_0 t_i}. &&
\end{flalign*}
(برای $n=0$ نیز چون $\sum_i \Delta_i=0$ در سیگنالِ دوره‌ای، نتیجه $g_0=0$ به‌دست می‌آید.)